\begin{pastexam}[label={px:gk3-109-ds-01}]{高考三級資訊處理 109 年 第一題}
（申論；15 分）

考慮數字 1 到 $n$，若將其順序重新排置，每個排列順序都稱作一個排列或置換（Permutation），例如
5 1 4 3 2 是 1 2 3 4 5 的一個排列。我們可以將一個數字 1 到 $n$ 的排列視為一個順序的映射 $P$，則前述
例子可表示為 $P(5) = 1$、$P(1) = 2$、$P(4) = 3$、$P(3) = 4$、$P(2) = 5$。當然，1 2 3 4 5 也是 1 2 3 4 5
的一個排列。在一個數字 1 到 $n$ 的排列 $P$ 中，若一對數字 $i$ 和 $j$，$1 \le i < j \le n$，
$P(j) < P(i)$，也就是在排列 $P$ 中較大的數字 $j$ 出現在較小的數字 $i$ 左邊（前面），我們稱此對數字
為反向（Inversion），而排列 $P$ 的反向數（Inversion number）則定義為排列 $P$ 中反向的總數量。請回答
下列問題：
\begin{itemize}[leftmargin=2em, labelsep=0.3em, itemsep=2pt]
\item[㈠] 數字 1 到 $n$ 的何種排列會有最大的反向數？最大反向數是多少？（5 分）
\item[㈡] 若給定一個數字 1 到 $n$ 的排列 $P$，請提出一個線性遞迴（Linear Recursive）的方式來算出
排列 $P$ 的反向數，並提供虛擬碼（Pseudo-code）與時間複雜度分析。（10 分）
\end{itemize}
\end{pastexam}

\begin{pastsoln}
注意題目的 $P$ 是「數字 $\to$ 位置」：$P(k)$ 是數字 $k$ 排在第幾個。 一對 $i < j$ 是反向，就是
$j$ 排在 $i$ 前面。 例子 $5\,1\,4\,3\,2$ 的反向有 $(1,5)$、$(4,5)$、$(3,5)$、$(2,5)$、$(3,4)$、
$(2,4)$、$(2,3)$，反向數是 $7$。

\textbf{㈠} 反向數最多是所有的數對都反向。 $1$ 到 $n$ 共有 $\binom{n}{2}$ 對，所以反向數
$\le \binom{n}{2} = n(n-1)/2$。 要每一對 $i < j$ 都有 $j$ 在 $i$ 前面，排列必須是由大到小
\[
n,\ n-1,\ \ldots,\ 2,\ 1 \qquad (\text{即 } P(k) = n + 1 - k),
\]
這時每一對都反向，反向數 $= n(n-1)/2$，達到上界；而且只有這一個排列能達到（任何別的排列都有某個
相鄰位置是小的在前，那一對不反向）。 例如 $n = 5$：$5\,4\,3\,2\,1$ 的反向數是 $10$。

\textbf{㈡ 線性遞迴。} 線性遞迴指每次呼叫最多只做一次遞迴呼叫。 把反向依「較大的那個數」分類：
令 $\mathit{Inv}(k)$ 是只看數字 $1..k$ 時的反向數。 數字 $1..k$ 之間的反向分成兩種：兩個數都
$\le k - 1$ 的，有 $\mathit{Inv}(k-1)$ 個；較大的那個是 $k$ 的，就是「比 $k$ 小、卻排在 $k$ 後面」
的數的個數，即 $c_k = |\{\, i < k : P(k) < P(i) \,\}|$。 所以
\[
\mathit{Inv}(1) = 0, \qquad \mathit{Inv}(k) = \mathit{Inv}(k-1) + c_k \quad (k \ge 2),
\]
答案是 $\mathit{Inv}(n)$。
\begin{pseudo}{以線性遞迴計算反向數}
\Function{Inversions}{$P, k$} \Comment{只看數字 $1..k$ 的反向數}
    \If{$k \le 1$}
        \State \Return $0$
    \EndIf
    \State $c \gets 0$
    \For{$i \gets 1$ \To $k - 1$}
        \If{$P(k) < P(i)$} \Comment{$i < k$ 而 $i$ 排在 $k$ 後面}
            \State $c \gets c + 1$
        \EndIf
    \EndFor
    \State \Return \Call{Inversions}{$P, k - 1$} $+\ c$ \Comment{只有一次遞迴呼叫}
\EndFunction
\end{pseudo}
呼叫 \proc{Inversions}($P, n$)。 \textbf{正確性}由上面的分類對 $k$ 歸納即得：每個反向 $(i, j)$ 恰好
在 $k = j$ 那一層被數到一次。

以 $5\,1\,4\,3\,2$ 為例（$P(1..5) = 2, 5, 4, 3, 1$）：$c_2 = 0$（$P(2) = 5$ 不比 $P(1) = 2$ 小）、
$c_3 = 1$（$P(3) = 4 < P(2) = 5$）、$c_4 = 2$（$P(4) = 3 < P(2), P(3)$）、$c_5 = 4$（$P(5) = 1$
比前四個都小），$\mathit{Inv}(5) = 0 + 1 + 2 + 4 = 7$。

\textbf{時間複雜度。} 第 $k$ 層的迴圈做 $k - 1$ 次比較，其他常數時間：
\[
T(k) = T(k-1) + \Theta(k), \quad T(1) = \Theta(1)
\;\Longrightarrow\;
T(n) = \Theta\Bigl(\sum_{k=1}^{n} k\Bigr) = \Theta(n^2).
\]
比較次數恰好是 $\sum_{k=2}^{n} (k-1) = n(n-1)/2$，和輸入無關。 遞迴深度 $n$，呼叫堆疊空間
$\Theta(n)$。

補充：若不限定線性遞迴，可以把排列寫成序列，用合併排序順便數反向（合併時右半的數先放出去，
就和左半剩下的每個數各構成一個反向），$T(n) = 2T(n/2) + \Theta(n) = \Theta(n \lg n)$，見資結書
第~17.5~節。
\end{pastsoln}
